\documentclass[a4paper,11pt]{article}
\usepackage{htlprotokoll}
\usepackage{graphicx}
\usepackage{pgfplots}
\pgfplotsset{compat=1.18}
\usepackage{lmodern} %Schriftart
\usepackage[ngerman]{babel}
\usepackage[T1]{fontenc}
\usepackage[utf8]{inputenc}
\usepackage{hyperref}
\usepackage{cleveref}
\crefname{figure}{Abb.}{Abb.}
\crefname{table}{Tab.}{Tab.}
\usepackage{float}


\labortext{zum Projekt}
\titel{RFID}
\klasse{4AHELS}
\datum{Sommersemester 2026 / 27.05.2026}
\lehrer{TILLICH und CHRA}

\personen
{Protokollführer}{Moritz Klein}
{Protokollführer}{Christoph Hackl}
{Mitarbeiter}{}
{Mitarbeiter}{}

\messobjekte{Messung A \\ Messung B}
\geraete{TTi 24-6.1 \\ Tektronix TBS 1052B}

\begin{document}
	
	% Deckblatt
	\protokolldeckblatt
	
	% Inhaltsverzeichnis
	\tableofcontents
	\newpage
	
	
	% Start mit Text
	\section*{Vorgaben}
	\begin{Center}
		$U_{e}=12V,U_{a}=9V,f_{r}=60kHz,P_{e}=P_{a}=3W$
	\end{Center}
	\section{Aufbau der NE555 Timerschaltung}
	Ziel dieser Aufgabe ist es mithilfe eines NE555-Timers ein Rechteckssignal für die weitere Schaltung zu erzeugen.
	\subsection{Schaltung}
	Die Schaltung \cref{fig:schaltung_ne555} besteht aus einem NE555 Timerbaustein
	\begin{figure}[h]
		\centering
		

		\caption{Schaltung der Timerschaltung}
		\label{fig:schaltung_ne555}
	\end{figure}
	
	\subsection{Messvorgang}
	\subsection{Messergebnis}
	\subsection{Interpretation}
	
	
	\section{Aufbau der Signalverzögerung und MOSFET-Endstufe}
	In dieser Aufgabe soll die bestehende Schaltung um eine Signalverzögerung, eine MOSFET-Treiberstufe und eine MOSFET-Gegentaktendstufe erweitert werden. Diese Erweiterungen dienen dazu, das PWM-Signal zu verzögern, die MOSFETs sicher anzusteuern und eine leistungsfähige Endstufe zu realisieren.
	
	\subsection{Schaltung}
	
	\begin{figure}[h]
		\centering
		

		\caption{Schaltung der Signalverzögerung und MOSFET-Endstufe}
		\label{fig:schaltung_endstufe}
	\end{figure}
	
	
	
	
	\noindent
	\textbf{Signalverzögerung:} Das PWM-Signal wird durch die NAND-Gatter in Kombination mit den RC-Gliedern verzögert. Dies sorgt für eine Totzeit, um gleichzeitiges Durchschalten der MOSFETs zu vermeiden.
	\\
	\textbf{MOSFET-Treiber:} Die verzögerten Signale werden durch mehrere CMOS-Inverter verstärkt. Diese dienen als Treiberstufe, um die MOSFETs mit ausreichend Spannung und Strom anzusteuern und schnelles Schalten zu ermöglichen.
	\\
	\textbf{MOSFET Gegentakt-Endstufe:} Die verstärkten Signale steuern die MOSFET-Endstufe, welche aus Q4 und Q5 besteht. Diese erzeugt ein leistungsstarkes PWM-Signal.
	
	\subsection{Messvorgang}
	\subsection{Messergebnis}
	\subsection{Interpretation}
	
	\section{Aufbau des Übertragers}
	In dieser Aufgabe soll die bestehende Schaltung um eine Signalverzögerung, eine MOSFET-Treiberstufe und eine MOSFET-Gegentaktendstufe erweitert werden. Diese Erweiterungen dienen dazu, das PWM-Signal zu verzögern, die MOSFETs sicher anzusteuern und eine leistungsfähige Endstufe zu realisieren.
	
	\subsection{Schaltung und Berechnung}
	Berechnet wird aus der Formel für eine Rahmenspule:
	\[
		L = N^2 \cdot d \cdot 10^{-6}
	\]
	wird zu:
	\[		
		d = \frac{L}{N^2 \cdot 10^{-6}} = \frac{150\,\mu H}{30^2 \cdot 10^{-6}} = 0,166666
	\]
	
	
	\begin{figure}[h]
		\centering
		
		
		\caption{Schaltung der Signalverzögerung und MOSFET-Endstufe}
		\label{fig:schaltung_übertrager}
	\end{figure}
	

	
	\subsection{Messvorgang}
	\subsection{Messergebnis}
	\subsection{Interpretation}
	
\end{document}
